
\subsubsection*{1.3.5}
\phantomsection
\label{sga3:VIIB:1.3.5}

Let \(k\) be a pseudocompact ring. A formal variety \(X\) over \(k\)
is said to be \emph{topologically flat} if its affine algebra \(A\) is
a \emph{projective} pseudocompact \(k\)-module, that is, if the
structure morphism
\[
  X\longrightarrow\operatorname{Spf}(k)
\]
is topologically flat.

Let us first note\footnote{\label{sga3:VIIB:nde:60}The editors added
the following lemma, used implicitly in the original; they also
introduced the numbering \SGARefLink{sga3:VIIB:1.3.5.A}{1.3.5.A}--\SGARefLink{sga3:VIIB:1.3.5.D}{1.3.5.D} for later references.} that
\SGARefLink{sga3:VIIB:0.2.2}{0.2.2} and \SGARefLink{sga3:VIIB:0.3.6}{0.3.6} imply the following result, analogous to
VII\textsubscript{A}, \SGARefLink{sga3:VIIA:3.1.1}{3.1.1}.

\medskip
\noindent\textbf{Lemma 1.3.5.A.}
\phantomsection
\label{sga3:VIIB:1.3.5.A}
{\itshape Suppose that \(k\) is Artinian. The functors
\[
  A\longmapsto A^\dagger=\operatorname{Hom}_c(A,k),
  \qquad
  C\longmapsto C^*=\operatorname{Hom}_k(C,k)
\]
define an anti-equivalence between the category of topologically flat
profinite \(k\)-algebras and the category of flat
\(k\)-coalgebras.}

Indeed, if \(A\) is a topologically flat profinite \(k\)-algebra,
then, by \SGARefLink{sga3:VIIB:0.3.6}{0.3.6}, there is an isomorphism of \(k\)-modules
\[
  A^\dagger\otimes_k A^\dagger
  \xrightarrow{\ \sim\ }
  \bigl(A\widehat{\otimes}A\bigr)^\dagger.
\]
Thus the multiplication
\[
  A\widehat{\otimes}A\longrightarrow A
\]
induces by duality a \(k\)-coalgebra structure on \(A^\dagger\). The
rest then follows from \SGARefLink{sga3:VIIB:0.2.2}{Proposition 0.2.2}.

Let us return to the case of an arbitrary pseudocompact ring \(k\).

\medskip
\noindent\textbf{Definition 1.3.5.B.}
\phantomsection
\label{sga3:VIIB:1.3.5.B}
To every formal \(k\)-variety \(X\) whose affine algebra \(A\) is a
projective pseudocompact \(k\)-module, one associates a flat
\(\mathbf O_k\)-coalgebra \(\mathbf H(X)\), defined as follows.

Denote by \(\mathbf H(X)\) the \(\mathbf O_k\)-module
\(\mathbf V_k^f(A)\), the ``dual of \(A\)''. It is a flat
\(\mathbf O_k\)-module, since the pseudocompact \(k\)-module
underlying \(A\) is projective \textup{(cf.\ \SGARefLink{sga3:VIIB:1.2.3.C}{1.2.3.C})}. Moreover, by
\SGARefLink{sga3:VIIB:0.3.6}{0.3.6}, there is an isomorphism
\[
  \mathbf V_k^f\bigl(A\widehat{\otimes}A\bigr)
  \simeq
  \mathbf V_k^f(A)\otimes\mathbf V_k^f(A).
\]
Consequently, the multiplication of \(A\) induces by transposition a
diagonal morphism
\[
  \mathbf H(X)=\mathbf V_k^f(A)
  \longrightarrow
  \mathbf V_k^f\bigl(A\widehat{\otimes}A\bigr)
  =\mathbf H(X)\otimes\mathbf H(X),
\]
which makes \(\mathbf H(X)\) a flat \(\mathbf O_k\)-coalgebra. We
shall say that \(\mathbf H(X)\) is the \emph{coalgebra of \(X\) over
\(\mathbf O_k\)}.

\medskip
\noindent\textbf{Definition 1.3.5.C.}
\phantomsection
\label{sga3:VIIB:1.3.5.C}
Conversely, to every \(\mathbf O_k\)-coalgebra \(\mathbf C\) one can
associate a \(k\)-functor \textup{(cf.\ \SGARefLink{sga3:VIIA:1.2.1}{1.2.1})}
\(\operatorname{Spf}^*(\mathbf C)\), defined as follows. For every
object \(B\) of \(\operatorname{Alf}/k\), set, with the notation of
VII\textsubscript{A} \SGARefLink{sga3:VIIA:3.1}{3.1},
\begin{equation}
\begin{aligned}
  \operatorname{Spf}^*(\mathbf C)(B)
  &=
  \operatorname{Hom}_{B\text{-coalg.}}\bigl(B,\mathbf C(B)\bigr)
  \\
  &=
  \bigl\{x\in\mathbf C(B)\mid
    \varepsilon_{\mathbf C(B)}(x)=1
    \ \text{and}\
    \Delta_{\mathbf C(B)}(x)=x\otimes x\bigr\}.
\end{aligned}
\tag{1}
\label{sga3:VIIB:1.3.5.C:1}
\end{equation}

Suppose in addition\footnote{\label{sga3:VIIB:nde:61}The editors
supplied the details that follow.} that the \(\mathbf O_k\)-module
underlying \(\mathbf C\) is admissible \textup{(cf.\ \SGARefLink{sga3:VIIA:1.2.1}{1.2.1})}, and put
\begin{equation}
  A
  =
  \Gamma^*(\mathbf C)
  =
  \varprojlim_{\mathfrak l}\mathbf C(k/\mathfrak l)^*.
\tag{2}
\label{sga3:VIIB:1.3.5.C:2}
\end{equation}
The algebra structures on the modules
\(\mathbf C(k/\mathfrak l)^*\) endow \(A\) with the structure of a
profinite \(k\)-algebra. For every object \(B\) of
\(\operatorname{Alf}/k\), one has
\begin{equation}
\begin{aligned}
  \operatorname{Hom}_{\operatorname{Vaf}/k}
    \bigl(\operatorname{Spf}(B),\operatorname{Spf}(A)\bigr)
  &=
  \operatorname{Hom}_{\operatorname{Alp}/k}(A,B)
  \\
  &=
  \operatorname{Hom}_{\operatorname{Alp}/B}
    \bigl(A\widehat{\otimes}B,B\bigr).
\end{aligned}
\tag{3}
\label{sga3:VIIB:1.3.5.C:3}
\end{equation}

Suppose finally that \(\mathbf C\) is a flat \(\mathbf O_k\)-module.
Then, by \SGARefLink{sga3:VIIB:1.2.3.E}{1.2.3.E},
\(A=\Gamma^*(\mathbf C)\) is a projective pseudocompact
\(k\)-module. Moreover, we saw in the proof of \emph{loc.\ cit.} that
if \(\mathfrak l_0\) is an open ideal of \(k\) contained in the kernel
of \(k\to B\), there are isomorphisms
\begin{equation}
\begin{aligned}
  A\widehat{\otimes}B
  &=
  \varprojlim_{\mathfrak l}
    \mathbf C(k/\mathfrak l)^*\widehat{\otimes}_kB
  \\
  &\simeq
  \operatorname{Hom}_{k/\mathfrak l_0}
    \bigl(\mathbf C(k/\mathfrak l_0),B\bigr)
  \\
  &\simeq
  \operatorname{Hom}_B\bigl(\mathbf C(B),B\bigr).
\end{aligned}
\tag{4}
\label{sga3:VIIB:1.3.5.C:4}
\end{equation}
We shall denote the term on the right by \(\mathbf C(B)^*\). Finally,
\SGARefLink{sga3:VIIB:1.3.5.A}{Lemma 1.3.5.A}, applied to the Artinian ring \(B\), gives a natural
isomorphism
\begin{equation}
  \operatorname{Hom}_{B\text{-coalg.}}\bigl(B,\mathbf C(B)\bigr)
  \xrightarrow{\ \sim\ }
  \operatorname{Hom}_{\operatorname{Alp}/B}
    \bigl(\mathbf C(B)^*,B\bigr).
\tag{5}
\label{sga3:VIIB:1.3.5.C:5}
\end{equation}
Combining \textup{\SGARefLink{sga3:VIIB:1.3.5.C:1}{(1)}}, \textup{\SGARefLink{sga3:VIIB:1.3.5.C:5}{(5)}}, \textup{\SGARefLink{sga3:VIIB:1.3.5.C:4}{(4)}}, \textup{\SGARefLink{sga3:VIIB:1.3.5.C:3}{(3)}},
and \textup{(2)}, we obtain, when \(\mathbf C\) is a flat
\(\mathbf O_k\)-coalgebra, an isomorphism of functors
\begin{equation}
  \operatorname{Spf}^*(\mathbf C)
  \simeq
  \operatorname{Spf}(A)
  =
  \operatorname{Spf}\bigl(\Gamma^*(\mathbf C)\bigr).
\tag{\(\star\)}
\label{sga3:VIIB:1.3.5.C:star}
\end{equation}

Consequently, if \(\mathcal A(X)\) denotes the affine algebra of a
formal \(k\)-variety \(X\), then, taking \SGARefLink{sga3:VIIB:1.2.3.E}{1.2.3.E} into account, we
obtain the following.

\medskip
\noindent\textbf{Proposition 1.3.5.D.}
\phantomsection
\label{sga3:VIIB:1.3.5.D}
{\itshape
\begin{enumerate}
  \item[\textup{(i)}]
  The functors
  \[
    X\longmapsto
    \mathbf H(X)=\mathbf V_k^f\bigl(\mathcal A(X)\bigr),
    \qquad
    \mathbf C\longmapsto
    \operatorname{Spf}^*(\mathbf C)
    =
    \operatorname{Spf}\bigl(\Gamma^*(\mathbf C)\bigr)
  \]
  define an equivalence between the category of topologically flat
  formal \(k\)-varieties and the category of flat
  \(\mathbf O_k\)-coalgebras.

  \item[\textup{(ii)}]
  Moreover, this equivalence ``commutes with base change'': if
  \(k\to K\) is a morphism of pseudocompact rings, then
  \(X\widehat{\otimes}_kK\) corresponds to
  \(\mathbf H(X)\otimes_kK\).
\end{enumerate}}

